\chapter{An exhaustive positive-$d$ source partition}
\label{component:pfb270019c56f}
\noindent\textbf{Component }\path{pfb270019c56f}\par
{\small\noindent Source: \path{CORE/prior/prior_core/inherited/positive_d_source_partition.tex}\par}
{\small\noindent SHA256: \path{fb270019c56f27e2c8a4ea33f4d82e8514970cd4c4e2a28aca98ffb85a4d502b}\par}
{\small\noindent Inventory use: Only displayed SP1-SP3 coefficient factorization and angular-balance proofs; missing terminal references are replaced by AU15-AU18 | Only actual SP2-SP3 equations and their bounded-source proof; no missing terminal reference is imported\par}
\medskip
\section{An exhaustive positive-$d$ source partition}
\label{pfb270019c56f:sec:positive-d-source-partition}

\noindent\textbf{Presentation scope note.} This manuscript imports this retained component only for the source partition (SP1)--(SP3) and its derivation. Historical all-tests terminal calls below are supplied by the complete replacement arguments named at their points of use; they are not additional imported premises. The original source file is unchanged.


Throughout this section the coercive canonical source is bounded and physical,
$d$ is a positive unit, and an actual finite noncommuting stationary
zero contact exists. Write $U=u_c^2,V=v_c^2$. For nonnegative pure
source quantities, $a\asymp b$ means that $a/b$ and $b/a$ are bounded
positive units. This notation is used only inside $K_c$.

\begin{lemma}[Two forbidden weak faces]
\label{pfb270019c56f:lem:positive-d-forbidden-faces}
Every positive-$d$ C2 source has $A$ a positive unit. In C1, at least
one of $m,c$ is a positive unit.
\end{lemma}
\begin{proof}
The following identities for the central determinant coefficient
follow by substituting into the displayed four-by-four parent:
\begin{align*}
 h\big|_{\lambda=\beta=A=0}
 &=-d(d+2n)(2D+2c+d)(2\alpha+d+2m),\\
 h\big|_{\lambda=n=m=c=0}
 &=-d(2D+d)(2\beta+d)(2A+2\alpha+d).
\end{align*}
Both are strictly negative when $d>0$. If $\bar A=0$ in C2,
the residue Laurent polynomial has only the terms
$h+q_{01}/U+q_{02}V/U+q_{12}V$. Substituting $U=t^2,V=1/t$
and sending the real number $t$ to infinity gives the negative limit
$h$, contradicting physicality. If $\bar m=\bar c=0$ in C1,
its residue Laurent polynomial is $h+q_{10}/V$; letting $V$ tend to
infinity gives the same contradiction. These are arguments over the
real residue source and involve no tensor division.
\end{proof}

\begin{proposition}[The positive-$d$ C2 source partition]
\label{pfb270019c56f:prop:c2-positive-d-partition}
The C2 source is covered by the existing small-$\alpha$ all-tests
theorem or its exact stabilizer image unless $\alpha,c,A$ are units.
In the latter case the actual contact satisfies
\[
 U\text{ bounded},\qquad V^{-1}\text{ bounded},\qquad
 U^2\asymp m,\qquad DV^2\asymp1.                 \tag{SP1}
\]
Consequently its four exhaustive source types are:
\begin{enumerate}
\item $m,D$ units: $u_c,v_c$ bounded nonzero in residue;
\item $m$ weak, $D$ a unit: $u_c$ infinitesimal and $v_c$ a bounded unit;
\item $m$ a unit, $D$ weak: $u_c$ a bounded unit and $v_c$ unbounded;
\item $m,D$ weak: $u_c$ infinitesimal and $v_c$ unbounded.
\end{enumerate}
The first three types have all-tests certificates from the preceding
modules. The fourth, with $n$ arbitrary, is the common reciprocal
type completed by \hyperref[component:p10686597065a]{RC1--RC25} (p.~\pageref{component:p10686597065a}), fitted to this source in \hyperref[component:pc0b2849f8d09]{AU16 and AU19} (p.~\pageref{component:pc0b2849f8d09}).
\end{proposition}
\begin{proof}
The stabilizer $(\lambda\ \beta)(m\ D)(c\ \alpha)$ exchanges
$c,\alpha$ and acts on the absolute half-angles by
$(|u|,|v|)\mapsto(1/|v|,1/|u|)$. At measurement level it exchanges
the parties and sends $(A_0,A_1;B_0,B_1)$ to
$(-B_0,B_1;-A_0,A_1)$, followed by local basis changes.
Thus either weak residue enters
the exhaustive CU/WA/WN union in \hyperref[component:pc0b2849f8d09]{AU15} (p.~\pageref{component:pc0b2849f8d09}). Lemma~\ref{pfb270019c56f:lem:positive-d-forbidden-faces}
gives the unit $A$. In the remaining branch write
\[
 q_{01}=ma_1,\quad q_{02}=mDa_2,\quad q_{10}=a_3,
 \quad q_{12}=Da_4,\quad q_{20}=na_5,\quad q_{21}=a_6,
\]
where all $a_i$ are bounded positive units. At the actual contact,
$L(U,V)=0$ and every noncentral term is nonnegative. Since $h$ is
bounded, $U,V^{-1},m/U,DV$ are bounded. Angular stationarity is
\[
 (q_{02}/U+q_{12})V^2=q_{10}+q_{20}U,
 \qquad U^2(q_{20}+q_{21}V)=mV(a_1+a_2DV).
\]
In the first equation the right side is a positive bounded unit,
and the left side is $DV^2(a_4+a_2m/U)$ with a positive-unit
parenthesis. Thus $DV^2\asymp1$. Dividing the second equation by
$V$ gives $U^2(a_6+na_5/V)=m(a_1+a_2DV)$; both parentheses
are positive bounded units, so $U^2\asymp m$. This proves (SP1)
and the exhaustive partition. Coercivity ensures $m,D>0$ as actual
source quantities; the same fact also follows from the two equations.

For the first type, the direct all-tests argument in
Lemma~\ref{pfb270019c56f:lem:positive-d-strong-dispatch} below applies. The second
has the natural pitchfork entry of \hyperref[component:pfb0cf454f653]{S} (p.~\pageref{component:pfb0cf454f653}) under the source fit and all-tests argument in \hyperref[component:pc0b2849f8d09]{AU7--AU9} (p.~\pageref{component:pc0b2849f8d09});
Lemma~\ref{pfb270019c56f:lem:positive-d-strong-dispatch} promotes that terminal
to all tests. The stabilizer exchanges $m,D$, so the third type is
the exact image of the second. No restriction on $n$ has been used.
\end{proof}

\begin{proposition}[The positive-$d$ C1 source partition]
\label{pfb270019c56f:prop:c1-positive-d-partition}
For C1, weak $\alpha$ or weak $A$ is already covered by an all-tests
theorem and its stabilizer image. Otherwise $\alpha,A$ are units.
After the C1 stabilizer, $c$ may be taken to be a unit. The remaining
exhaustive source alternatives are:
\begin{enumerate}
\item $m$ is a unit and $D$ or $\beta$ is a unit: the source has
an all-tests certificate;
\item $m$ is weak and $D$ is a unit: the source has an all-tests
pitchfork certificate;
\item $m$ is a unit and both $D,\beta$ are weak: the actual contact
has $u_c$ a bounded unit and $v_c$ unbounded;
\item $m,D$ are weak: the actual contact has $u_c$ infinitesimal
and $v_c$ unbounded, with $\beta$ arbitrary.
\end{enumerate}
The third type is now covered at every test by
\hyperref[component:p2cdf6197677b]{PP1--PP19} (p.~\pageref{component:p2cdf6197677b}), fitted to this source in \hyperref[component:pc0b2849f8d09]{AU18--AU19} (p.~\pageref{component:pc0b2849f8d09}). The fourth is covered at
every test when $\beta$ is a unit by
Corollary~\ref{pfb270019c56f:cor:c1-reciprocal-beta-all-tests} below. Thus only its
weak-$\beta$ subcase remains after these reductions; it is completed
by \hyperref[component:p98f24864b071]{CR1--CR39} (p.~\pageref{component:p98f24864b071}), fitted to this source in \hyperref[component:pc0b2849f8d09]{AU18--AU19} (p.~\pageref{component:pc0b2849f8d09}).
\end{proposition}
\begin{proof}
The C1 stabilizer
$(\lambda\ n)(m\ c)(A\ \alpha)(D\ \beta)$ exchanges $A,\alpha$.
Hence \hyperref[component:p0db05df21ad1]{WN1--WN44} (p.~\pageref{component:p0db05df21ad1}), together with the general entry in \hyperref[component:p3a3949869b9e]{WA1--WA10} (p.~\pageref{component:p3a3949869b9e}) and \hyperref[component:pc0b2849f8d09]{AU18} (p.~\pageref{component:pc0b2849f8d09}), cover either weak case. The forbidden-face lemma permits an orientation with $c$ a
unit in the remaining branch. Now
\[
 q_{01}=ma_1,\quad q_{02}=mDa_2,\quad q_{10}=a_3,
 \quad q_{12}=\beta m a_4+Da_5,\quad q_{21}=a_6,
 \quad q_{22}=\beta a_7,
\]
where all $a_i$ are positive bounded units. The contact equality
bounds each positive Laurent term, so $U,V^{-1},m/U,DV$ are bounded.
The $U$ stationarity equation reads
\[
 U^2(a_6+\beta a_7V)=m(a_1+a_2DV).               \tag{SP2}
\]
Its right parenthesis is a positive bounded unit, and its left
parenthesis is bounded below by a positive unit. Therefore weak
$m$ forces $U$ infinitesimal. If $m$ is a unit, boundedness of
$m/U$ forces $U$ to be a unit. The other stationarity equation is
\[
 a_3=V^2\{mDa_2/U+\beta ma_4+Da_5+\beta a_7U\}. \tag{SP3}
\]
If $m$ is a unit, so is $U$, and the braces are comparable with
$D+\beta$. Thus $V$ is a unit exactly when $D$ or $\beta$ is a
unit, and is unbounded exactly when both are weak. If $m$ is weak,
$U$ is infinitesimal. If $V$ is bounded, (SP2) gives $m/U$ weak,
and (SP3) forces $D$ to be a unit. Conversely a unit $D$ makes
$V$ bounded in (SP3), so the remaining branch has $V$ unbounded.
These alternatives exhaust all cases.

In the first type apply the stabilizer if necessary to arrange
$D$ a unit; then $d,c,A,D,\alpha,m$ are units and
Lemma~\ref{pfb270019c56f:lem:positive-d-strong-dispatch} applies. In the second,
\hyperref[component:pc0b2849f8d09]{AU17--AU18} (p.~\pageref{component:pc0b2849f8d09}) supplies the exact finite pitchfork entry for weak $m$ and unit $D$; \hyperref[component:pc0b2849f8d09]{AU7--AU9} (p.~\pageref{component:pc0b2849f8d09}) with the complete terminal \hyperref[component:pfb0cf454f653]{S} (p.~\pageref{component:pfb0cf454f653}) supplies the all-tests conclusion.
\end{proof}

\begin{lemma}[All-tests promotion for strong and pitchfork sources]
\label{pfb270019c56f:lem:positive-d-strong-dispatch}
Let C1 or C2 have $d,c,A,D,\alpha$ units. If $m$ is a unit, it has
a tensor certificate at every independent projective test. If $m$
is weak and its actual contact has $u_c$ infinitesimal and $v_c$
a bounded unit, the same conclusion follows from its finite-angle
pitchfork theorem.
\end{lemma}
\begin{proof}
All noncommuting residue tests are covered by
Theorem~\ref{p94dfbc8a5970:thm:all-noncommuting}. At the zero-zero angle corner,
a zero state is the $W$ ray because $\alpha$ is a unit. Normalize
$W=1$ and retain the entire nonnegative $m$ packet. Every remaining
monomial is in $(u,X,Y,Z)^2$. Its quadratic residue is
\[
 2(d+2\beta+2n)X^2+(4c+4D+2d)(u-Y)^2+4AY^2+4\alpha Z^2,
\]
which is positive definite in $(u,X,Y,Z)$. Finite polynomial
allocation therefore gives an exact strict-residue Gram; the
retained packet is restored. This works whether $m$ is weak or a
unit and includes every bounded projective state chart near the
unique zero ray.

At $v=0,u\ne0$, the zero-state classification excludes zeros because
$A,\alpha$ are units. At $u=0,v\ne0$, it excludes zeros if $m$ is
a unit. If $m$ is weak, the sole zero state is the $W$ ray. The
appropriate pitchfork theorem covers the matched remote residues
$v=\pm\bar v_c$; its printed positive quartic boundary factor
covers every unmatched finite remote residue by the strict
four-root Gram. Its projective remote endpoint is also strict.

All other projective angles use the exact homogeneous chart
identities in \hyperref[component:pc0b2849f8d09]{AU9} (p.~\pageref{component:pc0b2849f8d09}) and the complete single-infinity identities immediately following it.
Their hypotheses here are precisely the units $d,c,A,D$: at
$(\infty,0)$ use $(pX+Z,Y,W,vX)$, at $(0,\infty)$ use
$(X,Y,uW,Z)$, and at $(\infty,\infty)$ use the five-root
quadratic form on $(W,X,Y,p,q)$ after retaining the $\beta,\alpha,n$
packets.
Those displayed residue forms are positive definite independently
of the residues of $m,n,\alpha,\beta$. One infinite angle with the
other finite nonzero has no kernel by the same $c,A,D$ square
identities. This exhausts both projective angle lines and all state
charts; positive pure testing factors restore the original slack.
\end{proof}

\begin{corollary}[All tests for the strong-$\beta$ reciprocal source]
\label{pfb270019c56f:cor:c1-reciprocal-beta-all-tests}
Under the unit-$\beta$ subcase of \hyperref[component:p98f24864b071]{CR1--CR39} (p.~\pageref{component:p98f24864b071}) for weak $m,D$, fitted to the source in \hyperref[component:pc0b2849f8d09]{AU18--AU19} (p.~\pageref{component:pc0b2849f8d09}),
the C1 source has a tensor certificate at every independent projective
test, not only at its reciprocal testing corner.
\end{corollary}
\begin{proof}
That source fit gives $d,c,A,\alpha,\beta$ units, $m,D$ weak,
$\bar h=0$, and every state at the corner $(u,v)=(0,\infty)$.
Theorem~\ref{p94dfbc8a5970:thm:all-noncommuting} handles every noncommuting
residue test. At $(0,0)$ use the strict four-root Gram in
Lemma~\ref{pfb270019c56f:lem:positive-d-strong-dispatch}; its displayed residue
remains positive when $D$ is weak, since $c,A,\alpha,d$ are units.
At $v=0,u\ne0$ the units $A,\alpha$ exclude a zero state.

At $u=0$ and finite nonzero $v$, the only zero state is the $W$ ray.
Retain the whole $m$ packet and use $(u,X,Y,Z)$ as roots. The real
hard block is the positive $X$ mass together with
\[
 \{4(c+A)+2d\}(1+v^2)Y^2+4\alpha(vY-Z)^2,
\]
which is positive definite. The scalar angular Schur coefficient
is the pitchfork boundary expression
\[
 f(v)=\frac{8(1+v^2)}{t_7\det B(v)}
      (q_{10}+\bar h v^2+\bar q_{12}v^4)
      =\frac{8(1+v^2)q_{10}}{t_7\det B(v)}>0.
\]
Thus the exact four-root Gram has positive definite residue.
All divisions here describe a real residue calculation, not an
operation in the tensor ring.

At $(\infty,0)$ the existing four-root identity has residue
\[
 2dR^2+2d(W+V)^2+4cW^2+4AY^2\succ0,
\]
so it does not need a unit $D$. At $(\infty,\infty)$ the residue
state matrix is positive definite since $c,\beta,d$ are units.
For $u=\infty$ and finite nonzero $v$, the signed core together
with the $c,A,\beta$ packets contains
\[
 2d\{(W+vX)^2+(vY-Z)^2\}
 +4c(1+v^2)W^2+4A(1+v^2)Y^2
 +4\beta(1+v^2)Z^2\succ0.
\]
For $v=\infty$ and finite nonzero $u$, use instead
\[
 2d(1+u^2)(X^2+Y^2)+4c(uW-Y)^2
                              +4\beta(X+uZ)^2\succ0.
\]
These checks exhaust the projective angle charts. The stated
reciprocal theorem supplies the sole remaining double corner.
\end{proof}
